\documentclass[11pt]{article}
\usepackage[a4paper,margin=1in]{geometry}
\usepackage{amsmath,amssymb,amsthm,mathtools}
\usepackage{graphicx}
\usepackage[hidelinks]{hyperref}
\emergencystretch=3em

\newtheorem{theorem}{Theorem}
\newtheorem{lemma}[theorem]{Lemma}
\newcommand{\R}{\mathbb R}
\newcommand{\SB}{\mathrel{\perp_{\!SB}}}
\newcommand{\BJ}{\mathrel{\perp_{\!B}}}

\title{A Smooth Non-Euclidean Counterexample to the\\
Strong Birkhoff--James Basis Conjecture}
\author{A full counterexample to Conjecture 2.11 of arXiv:1407.5016}
\date{August 11, 2026}

\begin{document}
\maketitle

\begin{abstract}
We disprove the conjecture that a finite-dimensional real smooth normed space
of dimension at least three is Euclidean whenever every unit vector belongs to
a strongly orthonormal Hamel basis in the sense of Birkhoff--James.  An
explicit smooth, strictly convex, non-Euclidean example exists in every such
dimension.  It is the Euclidean direct sum of a non-Euclidean Radon plane and
a Euclidean space.  The required basis through each unit vector is given by a
closed formula.
\end{abstract}

\section{Source conjecture and answer}

Sain--Paul--Debnath ask in Conjecture 2.11 (page 6) whether the following
property characterizes inner-product spaces among finite-dimensional real
smooth normed spaces of dimension at least three: every point of the unit
sphere is contained in a strongly orthonormal Hamel basis in the sense of
Birkhoff--James.

\begin{center}
\includegraphics[width=.94\linewidth]{figures/open_problem_crop.png}
\end{center}

The answer is no, even if strict convexity is added.  We first construct an
explicit non-Euclidean Radon plane and then give counterexamples in all
dimensions $n\ge3$.

\section{Idea of Proof}

In a smooth space, $x\BJ y$ exactly when the unique norming functional at
$x$ annihilates $y$.  In a strictly convex space this already implies the
strict inequality defining $x\SB y$.  A Radon plane is a normed plane in
which Birkhoff--James orthogonality is symmetric.  Thus every unit vector
$a$ in such a plane has a unit partner $b$ with mutual orthogonality.

Now take the Euclidean direct sum of that plane and $\R^{n-2}$.  An arbitrary
unit vector has the form $(ra,sc)$ with $r^2+s^2=1$.  The vectors
$(ra,sc)$, $(-sa,rc)$, and $(b,0)$, followed by an orthonormal completion of
$c$ in the Euclidean summand, form the desired basis.  Their norming
functionals have the same two-dimensional rotation pattern, so every
off-diagonal pairing vanishes.

\section{An explicit smooth non-Euclidean Radon plane}

Fix $1<p<\infty$, $p\ne2$, and let $q=p/(p-1)$.  Define $N_p$ on $\R^2$ by
\begin{equation}\label{eq:mixed}
 N_p(x,y)=
 \begin{cases}
 (|x|^p+|y|^p)^{1/p},&xy\ge0,\\
 (|x|^q+|y|^q)^{1/q},&xy\le0.
 \end{cases}
\end{equation}
The definitions agree on the coordinate axes.
This mixed-exponent Radon construction is standard and is also recorded in
the source TeX accompanying the authors' related paper \cite{SPJ}; the
Euclidean-sum lifting below is the step that defeats the higher-dimensional
conjecture.  We include all details needed here.

\begin{lemma}\label{lem:radon}
The function $N_p$ is a smooth, strictly convex norm.  It is non-Euclidean,
and Birkhoff--James orthogonality in $E_p=(\R^2,N_p)$ is symmetric.
\end{lemma}

\begin{proof}
The upper boundary of the proposed unit ball is the graph
\[
 f(x)=
 \begin{cases}
 (1-|x|^q)^{1/q},&-1\le x\le0,\\
 (1-x^p)^{1/p},&0\le x\le1.
 \end{cases}
\]
Each piece is strictly concave, and their one-sided derivatives at zero are
both zero.  Hence $f$ is strictly concave.  The lower boundary is
$-f(-x)$, so the enclosed centrally symmetric set is convex and strictly
convex.  Formula \eqref{eq:mixed} is its Minkowski functional and therefore
a norm.  The two arcs have matching horizontal tangents at $(0,\pm1)$ and
matching vertical tangents at $(\pm1,0)$; elsewhere they are smooth.  Thus
every nonzero point has a unique norming functional.

Let $a=(a_1,a_2)$ be a unit vector off the axes and write $j_a$ for its
unique norming functional, identified with its coefficient vector.  If
$a_1a_2>0$, then
\[
 j_a=(\operatorname{sgn}(a_1)|a_1|^{p-1},
      \operatorname{sgn}(a_2)|a_2|^{p-1}).
\]
Put $b=(-(j_a)_2,(j_a)_1)$.  Its coordinates have opposite signs and,
using $(p-1)q=p$,
\[
 N_p(b)^q=|(j_a)_2|^q+|(j_a)_1|^q
          =|a_2|^p+|a_1|^p=1.
\]
Moreover its norming functional is $j_b=(-a_2,a_1)$, whence
$j_a(b)=j_b(a)=0$.  If $a_1a_2<0$, the same calculation interchanges $p$
and $q$ and uses $(q-1)p=q$.  The axis cases are immediate.  Since in a
smooth normed space $u\BJ v$ iff $j_u(v)=0$, orthogonality is symmetric.

Finally, $N_p(1,1)=2^{1/p}$ whereas $N_p(1,-1)=2^{1/q}$.  The parallelogram
identity would give
$2^{2/p}+2^{2/q}=4$, which, by strict convexity of
$t\mapsto2^t$ and $2/p+2/q=2$, holds only for $p=q=2$.
Thus $E_p$ is not Euclidean.
\end{proof}

\section{The counterexample in every dimension}

We use the following elementary smoothness observation.

\begin{lemma}\label{lem:strong}
Let $X$ be smooth and strictly convex, and let $j_x$ denote the unique
norming functional of $x\in S_X$.  Unit vectors $e_1,\ldots,e_n$ form a
strongly orthonormal Hamel basis in the sense of Birkhoff--James if
\[
 j_{e_i}(e_k)=\delta_{ik}\qquad(1\le i,k\le n).
\]
\end{lemma}

\begin{proof}
For $w$ in the span of all $e_k$ with $k\ne i$, the displayed identities
give $j_{e_i}(w)=0$, hence $e_i\BJ w$.  Therefore
$\|e_i+w\|\ge1$.  Equality for nonzero $w$ would put two distinct points
$e_i$ and $e_i+w$ on the unit sphere in the supporting hyperplane
$j_{e_i}=1$, contradicting strict convexity.  Hence
$\|e_i+w\|>1$ whenever $w\ne0$.
\end{proof}

\begin{theorem}\label{thm:main}
For every integer $n\ge3$ there is a smooth, strictly convex,
non-inner-product normed space $X_n$ in which every unit vector belongs to a
strongly orthonormal Hamel basis in the sense of Birkhoff--James.
Consequently Conjecture 2.11 of Sain--Paul--Debnath is false.
\end{theorem}

\begin{proof}
Let $H=\ell_2^{,n-2}$ and equip
\[
 X_n=E_p\oplus_2 H,
 \qquad \|(u,h)\|_X=(N_p(u)^2+\|h\|_2^2)^{1/2}.
\]
This norm is smooth and strictly convex because both summands are.  It is not
an inner-product norm, since its restriction to $E_p\oplus\{0\}$ is $N_p$.

Take $x=(u,h)\in S_{X_n}$ and put $r=N_p(u)$ and $s=\|h\|_2$, so
$r^2+s^2=1$.  If $r>0$, write $u=ra$ with $a\in S_{E_p}$; if $r=0$, choose
any $a\in S_{E_p}$.  By Lemma \ref{lem:radon}, choose $b\in S_{E_p}$ with
$j_a(b)=j_b(a)=0$.  If $s>0$, write $h=sc_1$ with $c_1\in S_H$; if $s=0$,
choose any $c_1\in S_H$.  Complete $c_1$ to a Euclidean orthonormal basis
$c_1,\ldots,c_{n-2}$ of $H$.

Define
\begin{align*}
 e_1&=(ra,sc_1)=x,&
 e_2&=(-sa,rc_1),&
 e_3&=(b,0),\\
 e_{k+2}&=(0,c_k) &&(2\le k\le n-2).
\end{align*}
The last line is absent for $n=3$.  These are unit vectors and form a basis:
$e_1,e_2$ are an orthogonal two-by-two change of coordinates in
$\operatorname{span}\{(a,0),(0,c_1)\}$, while the remaining displayed
vectors complete the two summands.

For a unit vector $(v,z)$ in the Euclidean sum, its norming functional is
\[
 j_{(v,z)}=\bigl(N_p(v)j_{v/N_p(v)},\,z\bigr),
\]
with first component zero when $v=0$.  Consequently
\[
 j_{e_1}=(rj_a,sc_1),\quad
 j_{e_2}=(-sj_a,rc_1),\quad
 j_{e_3}=(j_b,0),\quad
 j_{e_{k+2}}=(0,c_k),
\]
where vectors in $H$ are identified with their Euclidean functionals.
The relations $j_a(b)=j_b(a)=0$, $r^2+s^2=1$, and
$\langle c_i,c_k\rangle=\delta_{ik}$ now give
\[
 j_{e_i}(e_k)=\delta_{ik}
\]
for every pair of basis indices.  Lemma \ref{lem:strong} proves that this is
a strongly orthonormal Hamel basis containing the prescribed $x$.
\end{proof}

\section{Verification and novelty scope}

The accompanying deterministic script checks the mixed norm, its norming
functionals, the mutual Radon identities, and the full biorthogonality matrix
for many nondegenerate and edge-case vectors in dimensions three through
seven.  It also records the explicit failure of the parallelogram law.

A bounded search of the run indexes and exact terminology found the source
paper and subsequent general work on strongly orthonormal/Auerbach bases, but
no paper resolving Conjecture 2.11 by the Radon-plane Euclidean-sum
construction.  The novelty claim is limited to this explicit counterexample
and proof.

\begin{thebibliography}{9}
\bibitem{SPD}
D. Sain, K. Paul, and L. Debnath,
\emph{Characterization of inner product spaces}, arXiv:1407.5016 (2014),
Conjecture 2.11.

\bibitem{SPJ}
D. Sain, K. Paul, and K. Jha,
\emph{Strictly convex space: Strong orthogonality and conjugate diameters},
arXiv:1411.1464 (2014).
\end{thebibliography}

\end{document}
